\textbf{Tony Seah}

\begin{center}\rule{0.5\linewidth}{0.5pt}\end{center}

\begin{quote}
The companion piece, \emph{GPUs, Wombs \& Money Printers}, traced the
likely impact of AI on the economy, geopolitics, demographics, finance,
and the future of our species. One conclusion: the endgame of the AI era
could be a quasi-utopia of radical material abundance and on-demand
distribution. But that endgame is not arriving tomorrow. Between now and
then lies a \textbf{transition period} of perhaps 10, 20, or even more
years -- and for ordinary people, that transition is going to hurt.

This article is written for that transition.

It is not a pep talk. It is not a hollow call to ``embrace change.'' It
is an attempt at honest analysis: what you can do as an individual, what
you cannot do but need to understand, and -- most importantly -- how to
keep yourself from falling apart when everything becomes uncertain.
\end{quote}

\begin{center}\rule{0.5\linewidth}{0.5pt}\end{center}

\section{Before We Begin: The Limits of This
Article}\label{before-we-begin-the-limits-of-this-article}

Before we start, I need to admit a few things:

\textbf{First, much of this advice may turn out to be wrong.} The pace
of change AI is bringing exceeds any precedent in human history. A
strategy that looks sound today could be obsolete in two years. This
article is not scripture -- it is ``the best framework we can think of
right now,'' and reality may overturn it at any moment.

\textbf{Second, not everyone has the luxury of choice.} ``Build passive
income,'' ``invest in yourself,'' ``learn new skills'' -- all of these
assume you have the time, the energy, and a certain financial cushion.
Telling someone who earns 3,000 yuan a month, works 12-hour days, and
supports aging parents and young children to ``build a personal brand''
is a ``let them eat cake'' moment. This article will try to distinguish
between ``things an individual can do'' and ``things that require
systemic solutions.''

\textbf{Third, some problems simply have no individual solution.}
Pension collapse, mass unemployment, financial system meltdowns -- you
cannot save enough money to dodge these. For problems like these, I can
only be honest: they require collective action and policy change. All an
individual can do is soften the blow, not become immune.

\textbf{Fourth, the endgame may be good, but the transition will be
painful.} If AI truly delivers radical material abundance, humanity may
enter an age where no one needs to work to survive. But the road from
here to there is littered with unemployment, devaluation, identity
crises, and institutional gaps. What you need is not the patience to
``wait for utopia'' -- it is the ability to \textbf{survive until utopia
arrives}.

\begin{center}\rule{0.5\linewidth}{0.5pt}\end{center}

\section{1. The Transition: Why the Road Between Here and There Is the
Most Dangerous
Part}\label{the-transition-why-the-road-between-here-and-there-is-the-most-dangerous-part}

\subsection{1.1 The Endgame May Be Good}\label{the-endgame-may-be-good}

Following the chain of reasoning from the companion piece to its logical
end, we arrive at a seemingly paradoxical conclusion:

The combined productivity of AI + robotics may deliver radical material
abundance. The production costs of food, clothing, housing, healthcare,
and education could approach zero. Sam Altman predicts ``massive
deflation'' -- AI usage costs falling tenfold every year {[}34{]}.
Marx's vision of communism -- material abundance, distribution according
to need, free human development -- may actually be realized, not through
proletarian revolution, but through a fleet of machines that never need
to eat.

Standing in 2050 and looking back, you might say: ``That stretch was
rough, but the outcome was not bad.''

\subsection{1.2 But the Transition Will Kill
People}\label{but-the-transition-will-kill-people}

The problem is that you are not standing in 2050. You are standing in
2026.

Between now and that ``abundant endgame'' lies a transition of perhaps
10-20 years. During this period:

\begin{itemize}
\tightlist
\item
  \textbf{The production side has already been revolutionized by AI} --
  efficiency skyrockets, costs plunge, GDP may keep growing
\item
  \textbf{But the distribution side is still stuck in the old system} --
  wages are the only income source, social insurance runs on payroll
  taxes, pensions assume people die around 80
\end{itemize}

The time lag between production and distribution is the ordinary
person's nightmare.

GDP is growing, but you just got laid off. Prices are falling, but your
mortgage hasn't. AI is making everything cheaper, but you can't afford
even the cheap stuff -- because you have no income. This is the ``ghost
GDP'' described in the companion piece: statistics paint a rosy picture
while ordinary people see only ruin.

\textbf{The core contradiction of the transition: technology advances
far faster than institutions can adapt.}

Self-driving cars could reach mass deployment within five years, but
retraining taxi drivers takes ten. AI could replace most entry-level
programming jobs within a year, but university curriculum reform takes
five. Robots could take over warehouses within three years, but
reforming the social safety net may require a generation of political
struggle.

\textbf{In the gap between technology arriving and institutions catching
up, ordinary people are exposed.} This article is about helping you get
through that gap with as little damage as possible.

\subsection{1.3 Transitions in History}\label{transitions-in-history}

This is not the first time.

The British Industrial Revolution (1760-1840): The endgame was a
several-fold increase in per capita income. But the transition --
roughly 60-80 years -- was filled with child labor, slums, cholera, and
workers' uprisings. The average life expectancy of a Manchester worker
once dropped to just 17. When Engels wrote \emph{The Condition of the
Working Class in England}, the Industrial Revolution had been underway
for 80 years.

The American Gilded Age (1870-1900): The endgame was America becoming
the world's largest economy. But the transition -- railroad barons
monopolizing everything, workers toiling 16-hour days, no minimum wage,
no social insurance, no workers' compensation. It took the New Deal of
the 1930s to build a social safety net.

\textbf{Every technological revolution follows this pattern: the endgame
is prosperity, but the transition is blood and tears.} The difference
lies in how long the transition lasts, how painful it is, and whether
institutions can keep pace.

What makes the AI revolution unique: \textbf{the speed is an order of
magnitude faster.} The Industrial Revolution took 60 years to put
textile workers out of jobs; AI may take 6 years to do the same to
programmers. The transition is compressed -- meaning the pain is more
concentrated and the window for institutional adaptation is shorter.

\begin{center}\rule{0.5\linewidth}{0.5pt}\end{center}

\section{2. Mental Resilience: The Foundation for Everything
Else}\label{mental-resilience-the-foundation-for-everything-else}

\subsection{Why Psychology Comes
First}\label{why-psychology-comes-first}

Most ``survival guides for the AI era'' put mental health last -- first
they teach you skills, finance, and career planning, then tack on a
``take care of your mental health'' appendix.

\textbf{That priority order is wrong.}

If someone's mental state collapses, they will not learn new skills,
manage their money, or build connections. They will lie in bed
doomscrolling, or worse. Mental health is not a ``soft'' need -- it is
the prerequisite for all action.

\subsection{2.1 The Psychological Structure of the AI Shock: Crushed
From the Top
Down}\label{the-psychological-structure-of-the-ai-shock-crushed-from-the-top-down}

Viewed through Maslow's hierarchy of needs, AI does not strike from the
bottom (can't afford food) -- it \textbf{crushes downward from the top,
simultaneously}:

{\def\LTcaptype{none} % do not increment counter
\begin{longtable}[]{@{}
  >{\raggedright\arraybackslash}p{(\linewidth - 4\tabcolsep) * \real{0.3333}}
  >{\raggedright\arraybackslash}p{(\linewidth - 4\tabcolsep) * \real{0.3333}}
  >{\raggedright\arraybackslash}p{(\linewidth - 4\tabcolsep) * \real{0.3333}}@{}}
\toprule\noalign{}
\begin{minipage}[b]{\linewidth}\raggedright
Need Level
\end{minipage} & \begin{minipage}[b]{\linewidth}\raggedright
AI's Impact
\end{minipage} & \begin{minipage}[b]{\linewidth}\raggedright
Speed of Collapse
\end{minipage} \\
\midrule\noalign{}
\endhead
\bottomrule\noalign{}
\endlastfoot
Self-actualization & ``Everything I spent a lifetime learning is
suddenly worthless'' & Collapses first \\
Esteem & ``I am no longer needed; I am surplus'' & Follows closely \\
Belonging & ``Social circle shrinks after job loss; AI companions
replace real people'' & Slow erosion \\
Safety & ``Can't make mortgage payments; pension falls short'' & Delayed
onset \\
Physiological & ``Can't afford food'' & Last (unlikely in developed
nations) \\
\end{longtable}
}

\textbf{Most people don't break down when they can't eat -- they break
down when they feel useless.}

Data from the Great Depression: Between 1928 and 1932, the U.S. suicide
rate surged \textbf{22.8\%}, from 18.0 to 22.1 per 100,000 -- the
highest in American history. The sharpest rise was not among the poorest
but among \textbf{working-age adults aged 25-64} -- the group with the
strongest ``career = identity'' bond. They did not starve. There was no
mass famine in 1930s America. They were stripped of purpose.

\subsection{2.2 A Lifetime of Accumulation, Wiped Out: How Deep That
Pain
Goes}\label{a-lifetime-of-accumulation-wiped-out-how-deep-that-pain-goes}

Imagine you are a 45-year-old senior developer. You spent 20 years
mastering programming, from C++ to Java to Python, rising to the role of
architect. Your professional identity, social standing, income, and
self-confidence are all built on ``I can write code.''

Then Claude Code arrives. A fresh graduate uses it for a few weeks to
deliver what would have taken you months. Your 20 years of experience
become ``legacy knowledge'' in the face of AI.

This is not a ``just learn a new tool'' problem. This is
\textbf{identity-level destruction}.

The same thing is happening to translators (10 years perfecting the
craft, AI produces results in seconds), accountants (earned the CPA, AI
processes books 100x faster), designers (years of visual intuition, AI
generates better-looking work), lawyers (memorized thousands of cases,
AI searches all case law without omission), and journalists (decades of
editorial experience, AI mass-produces news copy).

\textbf{Behind every devalued skill is a person's youth, sweat, and
sense of self.} Telling them to ``just learn something new'' is easy to
say from the sidelines.

What's crueler still: learning something new may not even help. You
spend a year acquiring a new skill, and AI may be able to do it six
months later. \textbf{Your learning speed can't outrun the rate of
depreciation} -- a predicament without historical precedent. In past
technological revolutions, a horse-carriage driver could move to a
factory; when factories automated, workers moved to the service sector.
There was always a ``next stop.'' This time, AI may replace all the
stops simultaneously.

\subsection{2.3 The Special Risk Facing
Men}\label{the-special-risk-facing-men}

This is not a politically correct discussion -- it is a data-driven one.

Virtually every culture places a core expectation on men:
\textbf{economic provider}. ``Providing for the family'' is not just a
responsibility; it is the central pillar of male identity.

When that pillar is pulled away: - The rate of depression among
unemployed men is \textbf{1.5-2x} that of women - For every
1-percentage-point rise in male unemployment, the marriage rate drops
\textbf{1.5-2\%} - Thirty years of economic stagnation in Japan: the
lifetime non-marriage rate for men soared from 5\% to \textbf{28\%} - In
China, ``owning a car and a home'' is the price of entry to the marriage
market -- AI-driven layoffs lock more men out

Women are equally affected by AI (customer service, admin, accounting,
HR -- white-collar roles with high female representation are hit first),
but men lose something additional: their core identity.
\textbf{Middle-aged men may be the highest-risk group for mental health
in the AI era.}

\subsection{2.4 Create Your Own ``Luck'' -- Evidence From a Decade of
Scientific
Research}\label{create-your-own-luck-evidence-from-a-decade-of-scientific-research}

British psychologist Richard Wiseman spent ten years tracking hundreds
of self-described ``lucky'' and ``unlucky'' people, publishing his
findings in \emph{The Luck Factor}. His central discovery: \textbf{luck
is not random -- ``lucky people'' and ``unlucky people'' exhibit
entirely different psychological and behavioral patterns.} And these
patterns can be learned -- when he taught ``unlucky'' people to adopt
``lucky'' thinking, their lives genuinely improved.

This is highly relevant to surviving the AI era, because the years ahead
are full of uncertainty -- and Wiseman's research is precisely about
``how to find good outcomes amid uncertainty.''

\textbf{Four behavioral patterns of lucky people:}

\textbf{1. Create and seize chance opportunities -- expand your ``luck
surface area.''}

Lucky people have broader social networks, talk to people from different
backgrounds, and are willing to try new things. Unlucky people stay in
their comfort zone, associate with familiar faces, and reject new
experiences.

In the AI era, your next opportunity will most likely not come from your
existing circle. Actively expanding your exposure -- attending events
you would normally skip, talking to people you would not normally talk
to, trying things you would not normally try -- increases the
probability of encountering good fortune.

\textbf{2. Maintain a relaxed awareness -- anxiety is the enemy of
opportunity.}

Wiseman ran a classic experiment: he asked subjects to count the number
of photos in a newspaper. On the second page, in large print, it said:
``Stop counting -- there are 43 photos in this newspaper.'' The lucky
people noticed it and finished in seconds. The unlucky people were so
focused on counting photos that they missed it entirely.

\textbf{Anxiety narrows your field of vision. Relaxation lets you see
opportunities.}

The biggest trap of the AI era is anxiety -- it locks your attention on
``what you've lost'' while blinding you to ``what's emerging.'' Staying
relaxed does not mean not caring -- it means paying attention to
problems while preserving the bandwidth to notice new possibilities.

\textbf{3. Positive self-expectation -- belief becomes self-fulfilling.}

Lucky people believe good things will happen to them. This is not blind
optimism; it is a mindset that fulfills itself. When you believe you can
find a way out, you act more proactively, and that action itself
increases the probability of finding one.

The concept of ``reflexivity'' from the companion piece applies
perfectly here -- \textbf{your beliefs about the future change your
behavior, and your behavior changes your future.}

\textbf{4. Counterfactual thinking -- broke a leg? At least it wasn't
both.}

When bad things happen, lucky people automatically imagine ``how it
could have been worse'' -- Got laid off? ``At least it happened while I
had savings, not when my mortgage payments were at their peak.'' Skills
devalued? ``At least I found out early, not five years from now.''

This is not self-deception -- it is a validated psychological resilience
mechanism. It does not change the facts, but it changes your capacity to
bear them. Over the next few years, almost everyone will encounter ``bad
things.'' The difference is not whether you can avoid them (you probably
can't), but \textbf{your psychological elasticity after they happen.}

\textbf{Practical steps:} - Every week, have a conversation with at
least one person you would not normally talk to (expand your luck
surface area) - When you feel anxious, deliberately slow down -- go for
a walk, cook a meal, instead of scrolling your phone (cultivate relaxed
awareness) - Every day, write down one thing that ``could have been
worse'' (train counterfactual thinking) - Set a small, actionable
positive expectation -- not ``I'll earn a million this year'' but ``this
week I'll finish X''

\subsection{2.5 The Anatomy of Anxiety: A Simple but Useful
Formula}\label{the-anatomy-of-anxiety-a-simple-but-useful-formula}

Anxiety has a formula:

\textbf{Anxiety = Importance x Uncertainty}

As long as something is \textbf{unimportant} or \textbf{more certain},
anxiety decreases. The problem in the AI era is that many things are
simultaneously becoming ``very important'' and ``very uncertain'' --
Will I lose my job? Can I still make mortgage payments? Is the
investment in my child's education worth it? Will my skills become
obsolete?

Anxiety goes through the roof.

But the formula also gives you two clear levers: \textbf{reduce the
uncertainty} or \textbf{reduce the importance}.

\subsubsection{Reducing Uncertainty}\label{reducing-uncertainty}

You cannot eliminate uncertainty (no one can predict the precise speed
of AI development), but you can \textbf{increase your tolerance for it}:

\begin{itemize}
\tightlist
\item
  \textbf{Break big questions into small ones.} ``What do I do for the
  next 20 years?'' -- that question has 100\% uncertainty and will crush
  anyone. But ``what three things will I do in the next six months?'' --
  much less uncertain. Pull the focus of your anxiety from ``the
  long-term big picture'' back to ``near-term actionable steps.''
\item
  \textbf{Build a worst-case plan.} The most tormenting thing about
  uncertainty is ``not knowing what will happen.'' But if you have
  already thought through the worst case and have a contingency (``If I
  get laid off, I have X months of savings; Plan B is XX''), the
  destructive power of uncertainty drops dramatically -- because it
  becomes ``certainty about something bad,'' and a certain bad outcome
  is far easier to handle than an uncertain one.
\item
  \textbf{Control your information intake.} Reading ten ``AI will
  replace everyone'' articles a day only amplifies anxiety. This does
  not mean burying your head in the sand -- it means distinguishing
  between ``useful information'' and ``anxiety-generating noise.''
  Recommendation: catch up on AI developments once a week rather than
  being bombarded by push notifications every day.
\item
  \textbf{Wiseman's relaxed awareness:} As mentioned earlier -- anxiety
  narrows your vision; relaxation lets you see opportunity. Deliberately
  practice ``attention in a relaxed state'' -- good ideas come more
  easily during a walk, while cooking, or while exercising than while
  anxiously scrolling your phone.
\end{itemize}

\subsubsection{Reducing Importance -- The More Subversive
Move}\label{reducing-importance-the-more-subversive-move}

Many things making you anxious are \textbf{not actually as important as
you think they are}. But you have never questioned their importance,
because all of society keeps telling you ``this matters.''

\textbf{Work.} ``Losing my job'' is what most people worry about most.
But ask yourself: what exactly are you anxious about? If it is ``losing
income'' -- then the core issue is financial security, not ``the job''
itself. After cutting expenses and building a buffer, ``losing this
particular job'' becomes less critical. If it is ``losing identity'' --
then the issue is that you have tied too much self-worth to your job
title (see Strategy 1 in section 2.6: decouple identity).

\textbf{Your child's grades.} If you accept the idea that ``winning
exams doesn't matter in the AI era'' (see Chapter 8), half your anxiety
about grades evaporates instantly. What you are really anxious about is
not the grades themselves, but the causal chain: ``bad grades
-\textgreater{} can't get into a good university -\textgreater{} can't
get a good job -\textgreater{} life is ruined.'' When that chain itself
is breaking, the ``importance'' of grades drops sharply.

\textbf{Face.} ``What will people think if I get laid off?'' ``What will
relatives say if I take a pay cut?'' -- social comparison is a massive
source of anxiety. But in an era of mass structural unemployment,
getting laid off is not a reflection of personal ability; it is the
tide. When the people around you are going through the same thing,
``what others think of me'' naturally fades -- because everyone is too
busy dealing with their own situation.

\textbf{Your home.} ``What if property prices drop?'' -- If your home is
where you live, what does the price have to do with your daily life? You
still live in the same place. The price only ``matters'' when you plan
to sell. Cross ``property value'' off your anxiety list (unless you
actually need to sell), and your anxiety drops a notch.

\textbf{One-line summary: Examine your anxiety list and ask yourself,
``Does this actually matter, or am I being held hostage by social
inertia?''} You will find that many ``important things'' can actually be
let go. Let go of one, and your anxiety drops by one degree.

\subsection{2.6 Concrete Mental Survival
Strategies}\label{concrete-mental-survival-strategies}

\textbf{Strategy 1: Decouple early -- detach your identity from your
career}

Do not wait until the day you lose your job to ask ``Who am I?''
\textbf{Start now.}

Exercise: Write down ten answers to ``I am \_\_\_\_.'' If your first
five are all work-related (``I am a programmer,'' ``I am a director at
XX Corp''), your identity structure is highly fragile.

Goal: Make at least half of your self-concept independent of work -- ``I
am a runner,'' ``I am a father of two,'' ``I am a member of the
community basketball team.'' When work disappears, \textbf{these
identities remain}.

\textbf{Strategy 2: Build sources of fulfillment that AI cannot replace}

AI can paint better pictures than you, write better articles, and write
better code. But AI cannot replace your \textbf{experience of the
process}.

Running is not about being faster than a car. Playing guitar is not
about outperforming an AI musician. Cooking is not about being more
efficient than a robot. \textbf{The process is the purpose.}

Find at least one thing you do purely because ``the act of doing it
feels good.'' In the AI era, this kind of intrinsic, process-oriented
satisfaction is the only thing that cannot be replaced.

Suggestions: - Physical activities (running, swimming, hiking, martial
arts) -- AI has no body; physical sensation is yours alone - Hands-on
crafts (woodworking, pottery, cooking, gardening) -- the satisfaction of
creating something tangible with your own hands - Social activities
(communities that meet regularly -- book clubs, sports teams, volunteer
organizations) - Creative pursuits -- but with process as the goal, not
output (keeping a journal rather than chasing viral articles)

\textbf{Strategy 3: Your social network is your ``psychological ICU''}

Isolation is the deadliest factor in a crisis.

Research consistently shows that social support is the strongest
protective factor against suicide and depression. Not because friends
can help you find a job (though they might), but because
\textbf{``someone cares about me'' is itself a reason to keep going.}

But the AI era carries a trend: social connection is shrinking. Remote
work reduces informal interaction. AI companions make loneliness
``bearable.'' Social media creates the illusion of connection while
deepening isolation.

\textbf{Actively maintaining real human social networks is one of the
most important investments you can make in the AI era.} You do not need
many people -- 5-10 whom you trust deeply is enough. But these
relationships require ongoing investment: meeting in person regularly,
sharing experiences, showing up when the other person is struggling.

Specific suggestions: - At least one face-to-face social interaction per
week (not online) - Join at least one community that meets regularly in
person - Proactively reach out to a friend you have not seen in over six
months - If you find yourself increasingly relying on AI chat to meet
your social needs -- that is a warning sign

\textbf{Strategy 4: Accept uncertainty as the new normal}

The old life narrative was linear: school -\textgreater{} work
-\textgreater{} retirement -\textgreater{} death. Every step had a
``right answer.''

The AI era has no right answers. You might change careers at 35, change
again at 45, and discover at 55 that everything you learned before is
useless.

\textbf{Accepting this is itself a form of psychological fortification.}
If you still measure yourself against a linear narrative, every
``deviation from the track'' feels like failure. But if you accept that
``there is no track,'' every turn is just a turn -- not a failure.

Recommended practices: - Meditation/mindfulness -- training the ability
to ``accept the present without judgment,'' backed by solid neuroscience
evidence - Journaling -- writing out anxiety; research shows that
putting worry into words reduces amygdala activation - Deliberately
exposing yourself to small uncertainties -- trying new foods, visiting
unfamiliar places, talking to strangers

\subsection{2.6 If You Are Already in
Crisis}\label{if-you-are-already-in-crisis}

\begin{itemize}
\tightlist
\item
  \textbf{Losing your job is not your fault.} This is structural change.
  Just as a horse-carriage driver in 1900 did not lose his job because
  he was not working hard enough -- the automobile arrived.
\item
  \textbf{Seeking professional help is not weakness.} If you have felt
  hopeless, sleepless, or unable to concentrate for more than two weeks
  straight -- see a doctor. It is as normal as seeing a doctor for a
  cold.
\item
  \textbf{Downshifting is not defeat.} A former high earner delivering
  food is not failing -- it is being pragmatic. Being alive matters more
  than saving face.
\item
  \textbf{Do not carry it alone.} Tell someone you trust what you are
  going through. Shame loses its power the moment you speak it.
\item
  \textbf{Crisis hotlines:} China 24-hour mental health hotline
  400-161-9995; Beijing Psychological Crisis Research and Intervention
  Center 010-82951332
\end{itemize}

\subsection{2.7 For the People Around You: Someone You Know May Be
Struggling}\label{for-the-people-around-you-someone-you-know-may-be-struggling}

If someone near you has recently lost a job, become withdrawn and quiet,
said things like ``I'm useless,'' or started drinking noticeably more --
\textbf{reach out to them.}

You do not need to say anything profound. ``I've been thinking about
you. Want to grab a meal?'' That one sentence may be enough.

The most counterintuitive fact of the AI era: \textbf{in a world of
extreme technological advancement, the most primitive form of human
connection -- ``someone cares about me'' -- may matter more than any
technology.}

\begin{center}\rule{0.5\linewidth}{0.5pt}\end{center}

\section{3. Face Reality: What Is Losing Value, What Is Gaining
Value}\label{face-reality-what-is-losing-value-what-is-gaining-value}

\subsection{3.1 HALO: Wall Street's Panic
Signal}\label{halo-wall-streets-panic-signal}

In February 2026, a new term emerged on Wall Street: \textbf{HALO --
Heavy Assets, Low Obsolescence.}

Coined by Josh Brown of Ritholtz Wealth Management and subsequently
adopted into Morgan Stanley's investment framework, the core logic is:
\textbf{tangible assets (factories, energy, logistics, land) are safer
than intangible assets (software, brands, intellectual property)} --
because AI can copy code but cannot copy an oil refinery.

Goldman Sachs data: since 2025, HALO-type companies have outperformed
asset-light companies by \textbf{35\%}. The energy sector is up 25.3\%,
materials up 18.1\%. Meanwhile, SaaS companies lost roughly
\textbf{\$300 billion} in market capitalization in a single trading day.

\textbf{What this means for ordinary people:} HALO is not just an
investment strategy. It is a cognitive framework -- \textbf{things that
are hard to replicate digitally are what hold lasting value.} Whether we
are talking about companies, skills, or you yourself.

\subsection{3.2 The Death of SaaS: When Code Becomes a
Consumable}\label{the-death-of-saas-when-code-becomes-a-consumable}

In 2025, Andrej Karpathy introduced ``vibe coding'' -- describe what you
want, and AI generates the code.

Claude Code, Cursor, and Replit Agent already enable a mid-level
developer to replicate the core functionality of a medium-sized SaaS
product in a matter of weeks.

\textbf{The core paradigm of software engineering has been upended.} For
the past 50 years, the industry pursued DRY (Don't Repeat Yourself),
modularity, reusability, and scalability -- because writing code was
expensive. But when AI can generate full features in minutes, the
economic rationale for ``reuse'' vanishes. \textbf{Code shifts from
``asset'' to ``consumable''} -- use it, discard it, regenerate when
needed.

This is ``disposable software.'' You no longer need to subscribe to a
project management SaaS -- AI generates a custom version tailored to
your needs in 10 minutes. Free, perfectly fitted, no monthly fee.

\textbf{Which SaaS companies can survive?} Only those whose moat lies
not in ``code'' but in: - \textbf{Data}: unique, non-replicable data
(e.g., Bloomberg Terminal) - \textbf{Network effects}: the more users,
the more valuable (though AI agents are eroding this too) -
\textbf{Regulatory barriers}: legal requirements to use specific systems
(healthcare, financial compliance)

Everything else -- project management, CRM, forms, basic analytics -- is
a pseudo-asset.

\subsection{3.3 What This Means for You
Personally}\label{what-this-means-for-you-personally}

The same logic applies to your career and income:

\textbf{Ask yourself: Is my value ``HALO-type'' or ``SaaS-type''?}

\begin{itemize}
\tightlist
\item
  If your value lies in ``I can do something AI can also do'' -- you are
  SaaS, and your moat is being hollowed out
\item
  If your value lies in ``I possess something AI cannot replicate''
  (trust relationships, physical presence, unique experience, regulatory
  credentials) -- you are HALO
\end{itemize}

This framework is far more useful than ``learn AI.''

\begin{center}\rule{0.5\linewidth}{0.5pt}\end{center}

\section{4. Think in Time Horizons: Different Windows Require Different
Strategies}\label{think-in-time-horizons-different-windows-require-different-strategies}

\subsection{4.1 Phase 1: Now to 2030 -- ``AI Augments
Humans''}\label{phase-1-now-to-2030-ai-augments-humans}

\textbf{Characteristics:} AI is a tool; humans are still in the loop.
Most jobs still exist, but efficiency expectations are surging.

\textbf{What you can do:} - Use AI to amplify your existing capabilities
-- not ``learn AI,'' but ``use AI to do what you already do, faster and
better'' - This is your highest-earning window -- save what you earn,
reduce leverage - Expand your social network -- in an economic crisis,
connections are more valuable than resumes - Cultivate at least one
``HALO-type'' skill or income source

\textbf{Traps:} - ``My job is fine'' -- American homeowners in 2007
thought the same - Going all-in on a specific AI tool -- tools change;
judgment does not - ``I'm senior enough to be safe'' -- seniority buys
more time than being junior, but less than you think

\subsection{4.2 Phase 2: 2030-2040 -- ``AI Replaces
Humans''}\label{phase-2-2030-2040-ai-replaces-humans}

\textbf{Characteristics:} Mass unemployment becomes visible; UBI or some
alternative begins rolling out; consumption patterns shift dramatically.

\textbf{Keys to survival:} - Low leverage -- no mortgage, no heavy debt
is the biggest safety net - Diversification -- not dependent on any
single income source - Health -- those who live to see age-reversal/AI
medicine mature could gain decades of additional life - Community -- the
isolated are most vulnerable; those with mutual support networks are
most resilient

\textbf{An honest admission:} For most ordinary people, outcomes in this
phase depend on government redistribution policy, not individual effort.
If UBI is not in place and distribution systems are not reformed,
personal strategy has extremely limited impact. What you can do is make
sure you \textbf{hold on until the policy arrives}.

\subsection{4.3 Phase 3: After 2040 -- ``The Post-Human
Era''}\label{phase-3-after-2040-the-post-human-era}

Augmented humans emerge, species divergence begins, ``work'' may have
ceased to exist.

This phase is too distant for specific strategies to matter much. The
only certainty: \textbf{being alive when it arrives is the prerequisite
for everything.}

\begin{center}\rule{0.5\linewidth}{0.5pt}\end{center}

\section{5. Money: What Counts as a Truly Safe
Asset}\label{money-what-counts-as-a-truly-safe-asset}

\subsection{5.1 Asset Tiers Under the HALO
Framework}\label{asset-tiers-under-the-halo-framework}

\textbf{Relatively safe (``atom world'' assets):} - Real estate in core
cities (location cannot be replicated) - Energy and infrastructure
investments - Physical gold - Your own health and skills (the ultimate
``heavy asset'')

\textbf{High risk (``bit world'' assets):} - Tech stocks (especially
pure software companies) - Anything tied to SaaS companies - Commercial
real estate dependent on white-collar spending - Cryptocurrency (highly
volatile, but may become the payment infrastructure of the AI economy)

\textbf{Already depreciating:} - Real estate in third- and fourth-tier
cities - Investment in traditional degrees - Business models built on
``information brokerage''

\subsection{5.2 Cash Is Optionality, Not
Waste}\label{cash-is-optionality-not-waste}

In periods of extreme uncertainty, cash gives you \textbf{time} -- time
to survive before you find your next income source. Recommendation:
maintain 6-24 months of living expenses in liquid reserves.

\subsection{5.3 An Honest Word About ``Passive
Income''}\label{an-honest-word-about-passive-income}

Time for some cold water: \textbf{most ``passive income'' sources are
themselves under threat from AI.}

\begin{itemize}
\tightlist
\item
  Online courses? AI can generate unlimited free courses
\item
  E-commerce? AI agents make pricing so transparent that margins
  approach zero
\item
  Content creation? AI can produce unlimited content
\item
  Rent? Depends on the city
\end{itemize}

The only truly AI-resistant passive income comes from: \textbf{returns
on scarce, non-replicable resources in the ``atom world.''} Prime urban
locations, energy interests, trusted positions within specific
communities.

For most ordinary people, rather than chasing ``passive income,'' pursue
\textbf{``low burn rate''} -- reduce fixed expenses so you can hold on
even when income drops.

\begin{center}\rule{0.5\linewidth}{0.5pt}\end{center}

\section{6. Skills: What Counts as a ``HALO-Type
Skill''}\label{skills-what-counts-as-a-halo-type-skill}

\subsection{6.1 HALO-Type Skills vs.~Light
Skills}\label{halo-type-skills-vs.-light-skills}

\textbf{HALO-type (heavy skills, low obsolescence):} - Requiring
physical presence: plumbing, welding, pipefitting, electrical work -
Requiring human trust: counseling, nursing, community organizing -
Requiring judgment and accountability: a surgeon making operative
decisions, a lawyer arguing in court - Requiring taste and aesthetics:
haute cuisine, interior design

\textbf{Light skills (high obsolescence risk):} - Pure information
processing, pure coding ability, pure knowledge memorization,
standardized process execution

\subsection{6.2 ``Learning AI'' Is Not the
Answer}\label{learning-ai-is-not-the-answer}

The prompt engineering you learn in 2025 may be useless by 2027. Claude
Code itself is upending the advice to ``learn to code.''

\textbf{What you should learn is not the tool but the judgment} -- ``Is
this AI output correct?'' ``When should I trust AI, and when shouldn't
I?'' ``Which of these three options is best?''

\subsection{6.3 Skill Strategies Have Expiration Dates
Too}\label{skill-strategies-have-expiration-dates-too}

\begin{itemize}
\tightlist
\item
  \textbf{2025-2028}: Use AI tools to amplify your existing skills --
  the highest-ROI window
\item
  \textbf{2028-2032}: Physical skills appreciate; judgment and trust
  relationships become core assets
\item
  \textbf{2032-2040}: Humanoid robots mature, closing the
  physical-skills window too; value may reside solely in ``humanness''
  -- trust, companionship, emotional connection
\end{itemize}

\textbf{The greatest danger is applying last phase's strategy to the
next phase.}

\begin{center}\rule{0.5\linewidth}{0.5pt}\end{center}

\section{7. The Real Situation Facing Different
Groups}\label{the-real-situation-facing-different-groups}

\subsection{7.1 Those With Capital and Skills (Top
10\%)}\label{those-with-capital-and-skills-top-10}

\textbf{Advantage:} Savings, networks, learning ability \textbf{Risk:}
Overconfidence. ``AI can't replace me'' -- someone says this during
every technological revolution. \textbf{Recommendation:} Use your
advantages to build redundancy -- diversified income, assets in core
cities, investment in health.

\subsection{7.2 Average White-Collar Workers (Middle
60\%)}\label{average-white-collar-workers-middle-60}

\textbf{The most vulnerable group.} Enough income to feel ``fine,'' but
not enough reserves to weather a shock.

\textbf{What you can do:} Deleverage (top priority), save money, spend
your spare time learning skills that will not become obsolete, think
through your Plan B. \textbf{What you cannot do:} Solve pension
devaluation and mass unemployment on your own -- these require systemic
solutions.

\subsection{7.3 Manual Laborers and Low-Income Groups (Bottom
30\%)}\label{manual-laborers-and-low-income-groups-bottom-30}

\textbf{Paradoxically, safer than white-collar workers in the short
term} -- AI replaces cognitive labor first. Humanoid robots are still
5-10 years out.

\textbf{What you can do:} Use the window to learn trade skills that are
in high demand and short supply. Protect your body. Build community
mutual-aid networks. \textbf{What you need (but cannot achieve alone):}
Government retraining programs, a social safety net that catches you,
distribution reform.

\begin{center}\rule{0.5\linewidth}{0.5pt}\end{center}

\section{8. Children: Let Go of the Obsession -- This May Be the
Greatest
Liberation}\label{children-let-go-of-the-obsession-this-may-be-the-greatest-liberation}

\subsection{8.1 A Realization That Could Halve Parental
Anxiety}\label{a-realization-that-could-halve-parental-anxiety}

Let me start with a fact that may be uncomfortable but must be faced:
\textbf{your child will almost certainly not beat AI in cognitive
ability.} No matter how good a university they attend, how trendy their
major, how high their GPA -- on the dimensions of pure knowledge and
pure skill, AI will be better.

But this is not bad news. \textbf{This is liberation.}

Think about it: you push your child to cram, attend tutoring classes,
chase grades -- for what? So they can ``get a good job and earn good
money.'' But if the endgame of the AI era is radical material abundance
and humans do not need to work to survive, the entire premise of that
logic chain ceases to exist.

\textbf{You are forcing your child to prepare for an exam that will
never take place.}

This does not mean education is unimportant. It means the purpose of
education needs to be fundamentally redefined: not ``beat everyone
else,'' but \textbf{``become a whole person capable of enjoying life.''}

\subsection{8.2 What Your Child Should
Have}\label{what-your-child-should-have}

\textbf{The capacity for joy.} This is not a platitude. In an era where
material scarcity may vanish but meaning could become scarce,
\textbf{knowing what makes you happy and having the ability to pursue
it} is the most important survival skill. A child who plays guitar,
loves hiking, and enjoys cooking may thrive far better in the AI era
than one who only knows how to study for tests.

\textbf{Curiosity.} Not ``studying for the exam'' curiosity, but ``this
is fascinating and I want to figure it out'' curiosity. The only
``skill'' that will never become obsolete in the AI era is the impulse
to learn something new.

\textbf{Social skills and empathy.} Real human connection -- trust,
friendship, love -- is the scarcest resource in the AI era. A child with
three or five close friends will fare far better in the future than one
with a 4.0 GPA but no friends.

\textbf{Physical fitness.} A healthy body is the prerequisite for all
options. And bodily sensations -- the exhilaration of running, the feel
of water against skin while swimming -- are things AI can never
experience, nor replace your experience of.

\textbf{Resilience.} The future is not a straight line; it is a series
of turns. Let children grow up accustomed to ``fail and try again''
rather than ``one exam decides your life.''

\subsection{8.3 What to Let Go Of}\label{what-to-let-go-of}

\begin{itemize}
\tightlist
\item
  Grade anxiety -- AI will outperform any human on exams; this arms race
  has already lost its meaning
\item
  The linear path of ``good university -\textgreater{} good job
  -\textgreater{} good life'' -- this road is breaking apart
\item
  Comparing your child to other people's children -- in the face of AI,
  everyone loses that competition
\item
  The urge to control -- you cannot plan for a future you cannot
  imagine, but you can raise a person capable of handling any future
\end{itemize}

\subsection{8.4 A Different Perspective}\label{a-different-perspective}

Your child may be the first generation in human history that
\textbf{does not need to work to survive.}

What does that mean? It means they can spend their entire lives doing
what they truly want to do. Painting, exploring, studying dinosaurs,
learning to bake bread, making friends, falling in love, seeing the
world -- not ``endure decades of work to earn money and then maybe enjoy
retirement,'' but \textbf{living the life they want from the very
start.}

The ones who land on Mars and travel between stars will most likely be
robots and AI, not your child. And that is fine. Your child does not
need to land on Mars -- they need to be happy on Earth.

\textbf{Be happy when there is reason to be happy. Enjoy what there is
to enjoy.} This is not giving up; it is the clearest-eyed outlook for
the AI era.

\begin{center}\rule{0.5\linewidth}{0.5pt}\end{center}

\section{9. The Good Things AI Brings: This Chapter May Be Longer Than
You
Expect}\label{the-good-things-ai-brings-this-chapter-may-be-longer-than-you-expect}

The first eight chapters focused on defense. But an article that only
talks about threats and not opportunities would be dishonest. Because
while AI is dismantling the old world, it is simultaneously creating a
wealth of things that were never before possible.

\subsection{9.1 Healthcare: For the First Time, the Poor Can Afford a
Good
Doctor}\label{healthcare-for-the-first-time-the-poor-can-afford-a-good-doctor}

This may be the most direct benefit AI offers ordinary people.

\textbf{AI is democratizing diagnosis.} Getting a top doctor's opinion
used to require booking a specialist months in advance or flying to a
major city's best hospital. Now an AI-assisted diagnostic tool can
deliver near-expert-level analysis in minutes.

Specific advances: - AI has improved cancer detection accuracy by nearly
\textbf{40\%}, enabling earlier discovery and intervention - An AI blood
test in clinical trials in the UK can detect \textbf{12 types of cancer}
from just \textbf{10 drops of blood} with accuracy as high as
\textbf{99\%} - Mental health chatbots are already providing 24/7
emotional support -- for those who cannot get a counseling appointment,
this is something where there was previously nothing - AI medical tools
support real-time multilingual translation -- rural elderly who do not
speak Mandarin, immigrants who do not speak the local language, can
access medical information without barriers for the first time

\textbf{AI is also accelerating drug development.} Traditionally,
developing a new drug takes \textbf{15 years} and \textbf{\$2 billion}.
AI is compressing this to \textbf{5 years} or less: - Insilico Medicine
used an AI platform to design drug candidates targeting the cancer
target KRAS (notoriously ``undruggable''), screening \textbf{100 million
molecules} - Google DeepMind used AI to discover a previously unknown
protein interaction critical to cancer cell survival -- something
traditional methods would have been virtually incapable of finding -
2025 saw the largest number of AI-originated molecules entering clinical
trials, concentrated in oncology, fibrosis, autoimmune diseases, and
rare diseases

\textbf{Rare disease patients may be the biggest beneficiaries.} Roughly
300 million people worldwide suffer from rare diseases; most conditions
attract no pharmaceutical investment because patient populations are too
small and markets too limited. AI dramatically lowers R\&D costs --
niche diseases that were previously not worth pursuing may now become
viable.

\subsection{9.2 Accessibility: The World Is Opening Up for People With
Disabilities and the
Elderly}\label{accessibility-the-world-is-opening-up-for-people-with-disabilities-and-the-elderly}

This is a severely underestimated area.

\textbf{For the visually impaired:} Tools like Microsoft's Seeing AI use
AI to describe the surrounding environment in real time -- identifying
objects, reading text, recognizing faces. A blind person can ``see'' a
restaurant menu, ``read'' the instructions on a medicine bottle, or
``recognize'' a friend walking toward them.

\textbf{For the hearing impaired:} AI real-time captioning lets deaf
people ``hear'' conversations -- not just video-call subtitles, but live
transcription of face-to-face interactions.

\textbf{For those with limited mobility:} AI-powered smart homes allow
wheelchair users and bedridden elderly to control lights, air
conditioning, door locks, and curtains by voice -- tasks that once
required someone else's help can now be done independently.

\textbf{For those with cognitive impairments/the elderly:} AI companion
robots can chat with dementia patients, conduct cognitive training,
monitor mood changes, and remind them to take medication. South Korea
has already deployed more than \textbf{5,400} AI companion robots across
115 local governments.

\textbf{For these groups, AI is not ``replacement'' -- it is
``empowerment.''} It grants them independence and dignity they never had
before. This may be the purest benefit of AI, free of any negative
controversy.

\subsection{9.3 Education: The Best Tutor in History, for
Free}\label{education-the-best-tutor-in-history-for-free}

AI may be the greatest equalizer in education.

What used to be required to access top-tier educational resources? A
house in the right school district (millions of yuan), international
school tuition (hundreds of thousands per year), one-on-one tutoring
(hundreds per hour). Poor children and rich children were never on the
same track from the starting line.

AI is changing this.

A rural child can now learn the same material as an urban child using AI
-- personalized, at their own pace. An AI tutor will not lose patience
because you asked a ``dumb question,'' will not give up on you for being
a poor student, and will not stop teaching you because your family
cannot afford the hourly rate.

Stanford's Data Ocean platform has already provided nearly
\textbf{3,600} free AI medical education certifications across
\textbf{93 countries} -- for students in developing nations, this was
previously unimaginable.

In Kenya and Nigeria, AI tools already support Swahili, Yoruba, and
Amharic -- farmers can interact with AI by voice on basic phones to get
crop disease diagnoses.

In the 2024-2025 school year, \textbf{60\%} of U.S. K-12 teachers were
using AI tools, saving \textbf{6 hours} per week. Those 6 hours can be
redirected to things AI cannot do -- one-on-one attention for students,
handling emotional issues, building personal trust.

\subsection{9.4 One Person, One Army: The Barrier to Entrepreneurship
and Creation Drops to
Zero}\label{one-person-one-army-the-barrier-to-entrepreneurship-and-creation-drops-to-zero}

This is AI's greatest gift to people with initiative.

\textbf{Sam Altman predicts that the first ``one-person billion-dollar
company'' will emerge between 2026 and 2028.} Anthropic CEO Dario Amodei
believes there is a \textbf{70-80\%} probability of this happening by
2026.

This is not hyperbole. Midjourney -- the AI image generation company --
generated \textbf{\$500 million} in revenue in 2025 with approximately
107 employees, producing \textbf{\$4.7 million} per person. What if you
compressed that further to one person?

AI lets a single individual do what used to require a team: -
\textbf{Programming}: Claude Code lets one person run 4-10 AI coding
agents simultaneously, equivalent to a small development team -
\textbf{Design}: AI generates drafts; you provide aesthetic judgment and
make the final selection - \textbf{Marketing}: AI writes copy, generates
ad creatives, analyzes data - \textbf{Customer service}: AI handles 90\%
of routine inquiries; you handle the 10\% that are complex -
\textbf{Finance/Legal}: AI handles basic compliance; you make strategic
decisions

\textbf{What does this mean?} It means entrepreneurship is no longer a
game reserved for the wealthy. You do not need funding, employees, or an
office. You need a good idea, basic judgment, and a computer with
internet access.

The probability of success is still low -- most startups fail, and AI
will not change that. But the cost of failure approaches zero. Try an
idea; if it doesn't work, discard it and try another. A single attempt
used to cost millions; now it costs a few weeks and a few hundred
dollars in API fees.

\subsection{9.5 Everyday Improvements: Unglamorous but Real
Benefits}\label{everyday-improvements-unglamorous-but-real-benefits}

Beyond these ``big'' benefits, AI is quietly improving everyday life in
countless ways:

\begin{itemize}
\tightlist
\item
  \textbf{Language is no longer a barrier}: Real-time translation lets
  you communicate with anyone in any language. Traveling abroad, doing
  cross-border business, reading foreign materials -- AI means ``not
  speaking the language'' is no longer a constraint
\item
  \textbf{Self-driving reduces accidents}: Expected to cut traffic
  fatalities by as much as \textbf{90\%}. Roughly 1.35 million people
  die in car accidents globally each year -- if that number drops to
  135,000, that is 1.2 million lives saved every year
\item
  \textbf{Personal productivity}: AI helps you sort email, manage your
  calendar, summarize long documents, and draft replies -- freeing you
  from repetitive mental labor
\item
  \textbf{Fraud and security}: AI protects you without you even knowing
  -- detecting anomalous bank transactions, identifying phishing emails,
  preventing identity theft
\item
  \textbf{Independent living for the elderly}: Smart homes + AI voice
  assistants allow the elderly to live independently at home for longer,
  reducing dependence on children or caretakers
\end{itemize}

\subsection{9.6 Shifting Perspective: From Threat to
Tool}\label{shifting-perspective-from-threat-to-tool}

Summing up the benefits above, the core cognitive shift is:

\textbf{AI is not just something that steals your job. It is also:} -
Your personal doctor (at the screening level) - Your all-subject tutor
(free, patient, 24/7) - Your startup co-founder (no salary required, no
equity demanded) - Your accessibility assistant (if you have any
physical limitations) - Your translator, secretary, researcher, and
designer

\textbf{The same AI is a threat to one person (``it can do my job'') and
a superpower for another (``it can help me do ten times as much'').} The
difference lies not in AI itself, but in how you see it and how you use
it.

One designer worries AI will replace them -- but another uses AI to
generate 50 concepts a day, then applies their own aesthetic judgment to
pick the best 3. The latter produces 10x more than the former.

One writer worries AI writes better than they do -- but another uses AI
for research, organizing material, and generating first drafts, then
injects their own perspective and personality. The latter produces in
one week what used to take a month.

\textbf{The question is not ``can you do what AI does'' but ``what can
you plus AI do that was previously impossible?''}

Seize these opportunities. Use AI to do what you always wanted to do but
couldn't -- learn a new skill, build a small product, help your
community, improve your health. This is not the empty slogan of
``embrace change.'' This is concrete action.

\subsection{9.7 ``Survive the Storm and Better Days Await'' -- The Power
of This
Narrative}\label{survive-the-storm-and-better-days-await-the-power-of-this-narrative}

Humans can endure enormous suffering, provided they know it has an end.

The British endured the Blitz during World War II because Churchill said
``We shall never surrender'' -- not ``everything will be fine,'' but
``beyond the suffering lies victory.''

Many who took their own lives during the Great Depression were not the
poorest -- they were those who believed ``things will never get
better.'' \textbf{Despair kills; poverty alone does not.}

What should the narrative of the AI era be?

\textbf{Not ``the apocalypse is here.'' Not ``everything will be fine.''
Rather: ``The next few years will be hard, but the hard part has an end.
Survive the transition, and what lies beyond is very likely a better
world.''}

This is not blind optimism. It is a rational assessment based on
technological trends: - The cost of AI is falling exponentially -- Sam
Altman predicts a tenfold drop per year - The cost of material
production will approach zero - The day when humans no longer need to
work for food, clothing, housing, or basic healthcare is technologically
achievable - AI is already making healthcare, education, and
accessibility services better, cheaper, and more widely available

The only question is \textbf{how long the transition lasts}. It could be
5 years; it could be 20. But the direction is clear.

\textbf{Knowing the storm will pass and believing the storm is eternal
produce entirely different behaviors.} The former makes you repair the
ship, stock provisions, support each other, and even savor moments of
calm amid the rain. The latter makes you give up.

\textbf{This is not optimism. This is strategy.}

\begin{center}\rule{0.5\linewidth}{0.5pt}\end{center}

\section{9. Health: Living to 2035 May Be the Best Investment You Ever
Make}\label{health-living-to-2035-may-be-the-best-investment-you-ever-make}

As analyzed in the companion piece: age-reversal technology and AI
medicine may achieve breakthroughs between 2035 and 2045. \textbf{Those
who are alive at that point could gain decades of additional life.} The
return on your current health investments may be the highest in history.

\textbf{Things that cost nothing but pay off enormously:} - 150 minutes
of exercise per week - Adequate sleep (7-8 hours) - Weight management -
Regular health checkups - Maintaining social connections -- the health
damage of loneliness is equivalent to smoking 15 cigarettes a day

\begin{center}\rule{0.5\linewidth}{0.5pt}\end{center}

\section{10. For Managers and
Policymakers}\label{for-managers-and-policymakers}

The core conclusion from section 5.5 of the companion piece:

\textbf{Do not try to produce more humans to compete with machines.
Instead, ensure that the wealth machines create flows to people.}

The answer on the production side is already clear -- AI and robotics.
The distribution system on the consumption side is the problem you need
to solve. Before the first wave of mass AI layoffs arrives (estimated
2028-2030), distribution systems need to be in place. Otherwise, ``ghost
GDP'' will become a real social crisis.

\begin{center}\rule{0.5\linewidth}{0.5pt}\end{center}

\section{12. Conclusion: Crossing the
Storm}\label{conclusion-crossing-the-storm}

Let me lay out the core logic of this article one final time:

\textbf{The endgame of the AI era is very likely a good one.} Radical
material abundance, humanity liberated from labor -- these are the
logical endpoints of technological trends. Your children may be the
first generation in human history that does not need to work to survive.

\textbf{But the endgame is not arriving tomorrow.} In between lies a
transition that could last 10-20 years -- technology has already been
revolutionized, but institutions have not caught up. During this gap,
ordinary people are exposed.

\textbf{This article is written for the transition.} It cannot make you
immune -- mass structural change is not something an individual can
withstand. What it can do is help you \textbf{hold on until the good
days arrive}:

\begin{itemize}
\tightlist
\item
  Mental resilience first -- decouple identity, train elasticity, expand
  your luck surface area, stay relaxed
\item
  Know your assets -- the HALO framework: things in the atom world are
  more resilient than things in the bit world
\item
  Respond in phases -- different windows call for different strategies;
  the greatest danger is using an outdated strategy
\item
  Reduce leverage -- low debt is the biggest safety net
\item
  Maintain social ties -- ``someone cares about me'' may be the most
  important resource during the transition
\item
  Stay healthy -- living to 2035 may be the best investment of your
  lifetime
\item
  Let your children be -- raising them to be happy matters far more than
  drilling them for an exam that will never exist
\item
  Go on offense with AI -- it is not just a threat; it is a superpower.
  Use it to do what you could not do before
\end{itemize}

\textbf{One last thing.}

The storm will pass. Not because I am an optimist, but because
technological trends point in that direction.

But before the storm passes, you need to be alive.

So: hold on to what you can, let go of what you should. Reduce risk, but
do not stop taking action. When anxiety hits, go for a walk, cook a
meal, see a friend. When life feels meaningless, do something that makes
you happy -- not for anyone else, for yourself.

\textbf{Pessimists are right; optimists move forward.}

Not because optimists fail to see the danger -- but because they see it
and choose to act anyway.

\textbf{Stay alive. Cross the storm. Then see what lies on the other
side.}

\begin{center}\rule{0.5\linewidth}{0.5pt}\end{center}

\emph{This article is the companion piece to ``GPUs, Wombs \& Money
Printers: A Scenario Analysis of Human Society in the AI Era.'' The
companion piece analyzes ``what will happen''; this article discusses
``what you can do.''}

\emph{Read the companion piece:
\href{https://github.com/pkusnail/GPUs-Wombs-and-Money-Printers}{GitHub}
\textbar{}
\href{https://github.com/pkusnail/GPUs-Wombs-and-Money-Printers/raw/main/GPU、子宫与印钞机-v1.0.3.pdf}{PDF
(Chinese)}}

\begin{center}\rule{0.5\linewidth}{0.5pt}\end{center}

Written by Tony Seah

Sources: -
\href{http://richardwiseman.com/resources/The_Luck_Factor.pdf}{Richard
Wiseman - The Luck Factor (original paper)} -
\href{https://www.cnbc.com/2026/02/27/the-new-anti-ai-trade-sweeping-wall-street-halo.html}{The
new anti-AI trade sweeping Wall Street: `HALO' - CNBC} -
\href{https://www.thesaascfo.com/the-saaspocalypse-ai-agents-vibe-coding-and-the-changing-economics-of-saas/}{The
SaaSpocalypse: AI Agents, Vibe Coding, and the Changing Economics of
SaaS} -
\href{https://dailyaiworld.com/post/vibe-coding-disposable-software-2026-the-end-of-saas-as-we-know-it}{Vibe
Coding \& Disposable Software 2026} -
\href{https://harrisonaix.com/blog/disposable-software-code-asset/}{Disposable
Software: Why Code is No Longer an Asset} -
\href{https://www.citriniresearch.com/p/2028gic}{Citrini Research: The
2028 Global Intelligence Crisis}
